\documentclass[12pt]{article}
\usepackage[utf8]{inputenc}
\usepackage[T1]{fontenc}
\usepackage[margin=1in]{geometry}
\usepackage{setspace}
\usepackage{xcolor}
\usepackage[colorlinks=true, urlcolor=blue!60!black, linkcolor=blue!60!black]{hyperref}
\setlength{\parindent}{0em}
\setlength{\parskip}{0.6em}
\renewcommand{\baselinestretch}{1.0}
\begin{document}
\begin{center}
{\LARGE \textit{Lawrence v. Texas}, 539 U.S. 558 (2003)\par}
\vspace{0.5em}
{\large Case Brief --- Two-Page Edition\par}
\end{center}

\section{Memory Jogger}
\textit{Bowers} overruled: Texas's ban on private, consensual same-sex intimacy violates due process, and moral disapproval alone is no legitimate state interest (\href{https://www.courtlistener.com/opinion/130160/lawrence-v-texas/}{CourtListener}).

\section{Detailed Case Facts}
Deputies responding to a false weapons report entered John Lawrence's Houston apartment and found him with Tyron Garner in a private, consensual sexual act. Both were convicted under Texas's ``Homosexual Conduct'' statute, which criminalized same-sex but not opposite-sex sodomy.

\section{Procedural History}
Each man was fined \$200. The en banc Texas Court of Appeals affirmed, treating \textit{Bowers v. Hardwick} as controlling.

\section{Judicial Votes}
Reversed 6--3. Kennedy wrote for Stevens, Souter, Ginsburg, and Breyer; O'Connor concurred in the judgment; Scalia (with Rehnquist and Thomas) and Thomas dissented.

\section{Holding}
The Due Process Clause protects consenting adults' private, noncommercial intimate conduct in the home. Moral disapproval is not a legitimate interest, and \textit{Bowers} is overruled. The case did not involve minors, coercion, prostitution, or recognition of relationships.

\section{Analysis of Opinions}
\textbf{Majority:} \textit{Bowers} misframed the liberty interest; history is ``the starting point but not in all cases the ending point,'' and an ``emerging awareness'' counts. The law ``furthers no legitimate state interest which can justify its intrusion into the personal and private life of the individual.'' \textbf{O'Connor:} singling out same-sex conduct fails equal protection. \textbf{Scalia:} the ruling ``effectively decrees the end of all morals legislation.'' \textbf{Thomas:} a silly law, but no general right of privacy.

\section{Examples: Future Applications}
\textbf{Same-Side:} a ban on private sex between unmarried adults falls. \textbf{Opposite-Side:} prostitution and sex with minors remain punishable. \textbf{Fence-Sitter:} consensual adult incest or polygamy.

\section{Critique}
Originalists such as Nelson Lund and John McGinnis call the decision ``judicial hubris,'' and \textit{Dobbs} restored the history-and-tradition test. Laurence Tribe praised its ``double helix'' of liberty and equality, while Cass Sunstein noted its ambiguity about the standard of review.

\section{Key Quotations}
\begin{quote}
``Had those who drew and ratified the Due Process Clauses of the Fifth Amendment or the Fourteenth Amendment known the components of liberty in its manifold possibilities, they might have been more specific. They did not presume to have this insight. They knew times can blind us to certain truths and later generations can see that laws once thought necessary and proper in fact serve only to oppress. As the Constitution endures, persons in every generation can invoke its principles in their own search for greater freedom.''
\end{quote}
\noindent --- Justice Kennedy, Opinion of the Court, 539 U.S. at 578--79.

\end{document}
